% BEGIN REV09_FRICKE_JET
\subsection{The Fricke correspondence and differential recovery}
\label{corpus9:iv:fricke}

The level-two lattice trace and its Fricke completion
provide, after the exceptional construction, the
coordinates
\[
 t=t_2,\qquad s=\frac{4096}{t},\qquad
 T=t+s+24,\qquad t\in\mathbb C^\times.
\]
The involution exchanges \(t\) and \(s\). The two ordered
modular readings are
\[
 j_1=\frac{(t+256)^3}{t^2},\qquad
 j_2=\frac{(s+256)^3}{s^2}.
\]
These formulas use the uniformization already established in
\eqref{km:j-nivel2} and retain the choice of the oriented reading.

\begin{proposition}[Correspondence and discriminant]
\label{corpus9:iv:fricke-correspondencia}
The following identities hold:
\begin{align}
 j_1-j_2&=(s-t)(T+23),\\
 j_1+j_2&=T^2+T-7256,\\
 j_1j_2&=(T+248)^3.
 \label{corpus9:iv:fricke-simetricos}
\end{align}
The two values are roots of
\[
 X^2-(T^2+T-7256)X+(T+248)^3=0,
\]
whose discriminant is
\begin{equation}
 \Delta=(T-152)(T+104)(T+23)^2.
 \label{corpus9:iv:fricke-discriminante}
\end{equation}
\end{proposition}
\begin{proof}
Expanding the first reading with \(ts=4096\) gives
\(j_1=t+768+48s+s^2\); the second follows by exchanging
\(s,t\). Subtraction gives \((s-t)(t+s+47)\), and addition
reduces to \((t+s)^2+49(t+s)-6656\). Substitution of
\(t+s=T-24\) yields the first two identities. For the product,
\[
 j_1j_2
 =\frac{((t+256)(s+256))^3}{(ts)^2}
 =(T+248)^3.
\]
Finally,
\((s-t)^2=(T-24)^2-16384=(T-152)(T+104)\).
Multiplication by \((T+23)^2\) proves
\eqref{corpus9:iv:fricke-discriminante}.
\end{proof}

\begin{proposition}[Recovery of orientation from values and the first jet]
\label{corpus9:iv:fricke-jet}
For \(T\ne-23\), the ordered pair \((T,j_1)\) recovers
\begin{equation}
 t=\frac{T^2-j_1-3904}{T+23}.
 \label{corpus9:iv:fricke-inversa}
\end{equation}
At \(T=-23\), the two parameters
\[
 t_\pm=\frac{-47\pm45i\sqrt7}{2}
\]
share \((T,j_1)=(-23,-3375)\). The first jet
\(p=dj_1/dT\), evaluated on the branch parametrized by \(t\),
separates them and recovers
\begin{equation}
 p_\pm=\frac{-45\mp45i\sqrt7}{2},\qquad
 s=p-1,\qquad t=-46-p.
 \label{corpus9:iv:fricke-jet-inversa}
\end{equation}
\end{proposition}
\begin{proof}
The identity \(s^2=(T-24)s-4096\) gives
\[
 j_1=(T+23)s+T-3352
     =T^2-3904-(T+23)t.
\]
For \(T\ne-23\), this yields
\eqref{corpus9:iv:fricke-inversa}. For \(T=-23\),
the conditions \(t+s=-47\), \(ts=4096\) give
\(t^2+47t+4096=0\), with discriminant \(-45^2\cdot7\).
The same identity gives \(j_1=-3375\) at both roots.

The derivative \(dT/dt=1-4096/t^2\) vanishes only at
\(t=\pm64\). Neither exceptional root has that value,
so \(T\) is a local coordinate on both branches.
Differentiating the first expression for \(j_1\),
\[
 \frac{dj_1}{dt}
  =(s+1)\frac{dT}{dt}+(T+23)\frac{ds}{dt}.
\]
At \(T=-23\), this gives \(p=s+1=-46-t\), proving
\eqref{corpus9:iv:fricke-jet-inversa}, with opposite signs
for the paired indices \(t_\pm,p_\pm\).
\end{proof}

The fixed points of Fricke, \(t=\pm64\), yield respectively
\((T,j)=(152,8000)\) and \((-104,1728)\). The coincidence at
\(T=-23\) instead corresponds to two distinct points of the
parametrization. The field of definition follows from the equation:
\[
 \mathbb Q(t_\pm)=\mathbb Q(\sqrt{-7}),\qquad
 \left(\frac{2t+47}{45}\right)^2=-7.
\]
On this fibre, Fricke agrees with quadratic conjugation,
with trace \(-47\) and norm \(4096\). The jet retains this field,
since \((2p+45)^2=-45^2\cdot7\). Reconstruction is exact
within this modular correspondence: the differential datum separates
the two orientations that share their scalar values.
% END REV09_FRICKE_JET
